\uuid{FdTF}
\exo7id{684}
\titre{exo7 684}
\auteur{ridde}
\organisation{exo7}
\datecreate{1999-11-01}
\isIndication{false}
\isCorrection{false}
\chapitre{Continuité, limite et étude de fonctions réelles}
\sousChapitre{Continuité : pratique}
\cpge{mpsi-mp2i.nombres-reels-suites.i, pcsi.nombres-reels-suites.g, ptsi.nombres-reels-suites.g, mpsi-mp2i.nombres-reels-suites.d, pcsi.nombres-reels-suites.c, ptsi.nombres-reels-suites.c, mpsi-mp2i.nombres-reels-suites.e, pcsi.nombres-reels-suites.d, ptsi.nombres-reels-suites.d, bcpst1.suites-reelles, bcpst1.limites-continuite.c}
\module{Analyse}
\niveau{L1}
\difficulte{}

\contenu{
\question{
En \'etudiant la suite
 $u_0 \in \Rr$ et $u_{n  + 1} = \cos (u_n)$, d\'eterminer une valeur approch\'ee \`a $10^{-5}$
 pr\`es de l'unique r\'eel solution de $\cos (x) = x$.
}
}
